\section{Hecke特征值的代数性}\label{sec:modular-2-2}

\subsection{Jacobi簇、同调与Abel定理}\label{sec:modular-2-2-1}

设 $X$ 是亏格为 $g$ 的紧Riemann曲面.若 $\omega\in\Omega^1(X)$ 是全纯微分形式，就可以沿任意连续路径 $c:[0,1]\longrightarrow X$ 对它积分.由留数定理及 $\omega$ 的全纯性，若 $c'$ 与 $c$ 同伦，则 $\int_{c'}\omega=\int_c\omega$.固定一点 $x_0\in X$.限制到以 $x_0$ 为基点的闭路，就得到群同态
\begin{align*}
\pi_1(X,x_0)_{\mathrm{ab}}=H_1(X,\mathbb Z)\longrightarrow\Omega^1(X)^*.
\end{align*}
用一点代数拓扑可以看出，Abel群 $H_1(X,\mathbb Z)$ 是秩为 $2g$ 的自由群，由环绕 $X$ 的 $g$ 个``洞''的纵向闭路 $c_1,\ldots,c_g$ 与横向闭路 $c'_1,\ldots,c'_g$ 生成.再用一点Hodge理论可以证明，族
\begin{align*}
\left\{\int_{c_i},\ i=1,\ldots,g\right\}\cup\left\{\int_{c'_i},\ i=1,\ldots,g\right\}
\end{align*}
是 $\Omega^1(X)^*$ 的一组 $\mathbb R$ 基.特别地，$H^1(X,\mathbb Z)$ 是复向量空间 $\Omega^1(X)^*$ 中的余紧格\footnote{译注：原文此处印作 $H^1(X,\mathbb Z)$；上下文中的积分映射与下式使用的是同调群 $H_1(X,\mathbb Z)$.}，商
\begin{align*}
J(X):=\Omega^1(X)^*/H_1(X,\mathbb Z)
\end{align*}
自然成为维数为 $g$ 的紧复流形，且带有解析Abel群结构（称为``复环面''）.



现在，若 $c$ 是从 $x_0$ 到 $x$ 的路径，则元素 $\int_c\in\Omega^1(X)^*$ 在 $J(X)$ 中的像只依赖于 $x$，因而可以简记为 $\int_{x_0}^x$.于是考虑映射
\begin{align*}
\operatorname{Div}^0(X)\longrightarrow J(X),\qquad
\sum_x n_x[x]\longmapsto\sum_x n_x\int_{x_0}^x.
\end{align*}
Abel--Jacobi定理断言，此映射诱导群同构\footnote{这证实了前面用过的事实：群 $\operatorname{Pic}^0(X)$ 是可除群.}
\begin{align*}
\operatorname{Pic}^0(X)\xrightarrow{\sim}J(X).
\end{align*}
下面将证明，当 $X=X_0(N)$ 时，所有这些对象都带有Hecke算子 $T_p$ 的自然作用，且此作用与Abel--Jacobi映射及同构 $\mathcal S_2(\Gamma_0(N))\simeq\Omega^1(\Gamma_0(N))$ 相容\footnote{译注：原文此处的右端印作 $\Omega^1(\Gamma_0(N))$；按上下文应为 $\Omega^1(X_0(N))$，此处保留原记号.}.



\subsection{函子性}\label{sec:modular-2-2-2}

设 $\pi:X\longrightarrow Y$ 是紧Riemann曲面之间的非常值态射.它诱导两个同态
\begin{align*}
\pi_*:\operatorname{Pic}^0(X)\longrightarrow\operatorname{Pic}^0(Y),\qquad
\pi^*:\operatorname{Pic}^0(Y)\longrightarrow\operatorname{Pic}^0(X),
\end{align*}
分别由 $\pi_*([x]):=[\pi(x)]$ 与 $\pi^*([y])=\sum_{x\in\pi^{-1}(y)}e_x[x]$ 定义，其中 $e_x$ 是 $\pi$ 在 $x$ 处的分歧指数.这两个在 $\operatorname{Div}$ 上的定义显然保持 $\operatorname{Div}^0$.还须说明它们保持主除子.为此，注意到 $\pi^*(\operatorname{div}(f))=\operatorname{div}(f\circ\pi)$，而 $\pi_*(\operatorname{div}(f))=\operatorname{div}(N(f))$，其中 $N$ 是扩张 $\mathbb C(X)/\mathbb C(Y)$ 的范数（具体地，$N(f)(y)=\prod_{x\in\pi^{-1}(y)}f(x)^{e_x}$）.



它还诱导两个 $\mathbb C$ 线性映射
\begin{align*}
\pi_*:\Omega^1(X)\longrightarrow\Omega^1(Y),\qquad
\pi^*:\Omega^1(Y)\longrightarrow\Omega^1(X).
\end{align*}
\begingroup\emergencystretch=3em
我们已经遇到过拉回 $\pi^*$，在局部坐标 $u=\pi(z)$ 下，其定义就是 $\pi^*(f(u)\,du)=f(\pi(z))\pi'(z)\,dz$.前推 $\pi_*$ 的定义稍微细致一些.先在 $Y$ 的开子集 $Y'$ 上定义 $\pi_*(\omega)$，其中 $\pi$ 在 $Y'$ 上方无分歧.对 $y\in Y'$，选取 $y$ 的邻域 $V$，使得 $\pi^{-1}(V)=\bigsqcup_{x\in\pi^{-1}(y)}U_x$ 且 $\pi|_{U_x}:U_x\xrightarrow{\sim}V$.置
\par\endgroup
\begin{align*}
\pi_*(\omega)|_V:=\sum_x\bigl((\pi|_{U_x})^{-1}\bigr)^*(\omega|_{U_x}).
\end{align*}
接着须证明它可以全纯延拓到 $Y$.在分歧点附近考察，可归结为 $\pi$ 在以 $0$ 为中心的去心圆盘 $D^*$ 上具有形式 $z\mapsto u=z^e$ 的情形.计算得到
\begin{align*}
\left.\pi_*\left(\sum_{n\in\mathbb N}a_nz^n\,dz\right)\right|_{D^*}
=\sum_{m\in\mathbb N}a_{e(m+1)-1}u^m\,du,
\end{align*}
右端显然可以全纯延拓到 $0$.



最后，还有映射
\begin{align*}
\pi_*:H_1(X,\mathbb Z)\longrightarrow H_1(Y,\mathbb Z),\qquad
\pi^*:H_1(Y,\mathbb Z)\longrightarrow H_1(X,\mathbb Z).
\end{align*}
映射 $\pi_*$ 就是对闭路取复合 $c\mapsto\pi\circ c$；与前面一样，$\pi^*$ 的定义较为细致.从 $Y$ 上以 $y$ 为基点的闭路 $c$ 出发，先将它变形，使之避开分歧轨迹.因为 $\pi$ 是局部同构，对每个 $x\in\pi^{-1}(y)$，都存在唯一提升 $c$ 的路径 $c_x$，满足 $c_x(0)=x$.这样便得到纤维 $\pi^{-1}(x)$ 的置换 $x\mapsto c_x(1)$\footnote{译注：原文将此纤维印作 $\pi^{-1}(x)$；这里的基点为 $y$，依上下文应为 $\pi^{-1}(y)$.}.若 $x_0\mapsto x_1\mapsto\cdots\mapsto x_r$ 是此置换的一个轨道，则各 $c_{x_i}$ 的首尾拼接构成 $X$ 上的一个闭路.定义 $\pi^*(c)$ 为各轨道所对应闭路之和的像.由定义可验证下面的引理.

\begin{lemma*}

对 $Y$ 上的闭路 $c$ 及 $\omega\in\Omega^1(X)$，有 $\int_{\pi^*c}\omega=\int_c\pi_*\omega$.对 $X$ 上的闭路 $c$ 及 $\omega\in\Omega^1(Y)$，有 $\int_{\pi_*c}\omega=\int_c\pi^*\omega$.

\end{lemma*}

此引理说明，$\Omega^1$ 上 $\pi^*$ 与 $\pi_*$ 各自的对偶映射诱导态射
\begin{align*}
\pi_*:J(X)\longrightarrow J(Y),\qquad\pi^*:J(Y)\longrightarrow J(X).
\end{align*}
同样，由定义还可验证：

\begin{lemma*}

Abel--Jacobi映射 $\operatorname{Pic}^0\longrightarrow J$ 与 $\pi_*$ 和 $\pi^*$ 相容.

\end{lemma*}

最后，我们关心的是下述情形.

\begin{lemma*}

设 $X=X(\Gamma)$、$Y=X(\Gamma')$，其中 $\Gamma\subset\Gamma'$，并分别记 $\pi_\Gamma$、$\pi_{\Gamma'}$ 为从 $\mathbb H^*$ 到 $X(\Gamma)$、$X(\Gamma')$ 的投影.则下面两个图交换：
\[
\begin{tikzcd}[column sep=large,row sep=large]
\Omega^1(X(\Gamma)) \arrow[r,"\sim","\pi_\Gamma^*"'] & \mathcal S_2(\Gamma) \\
\Omega^1(X(\Gamma')) \arrow[u,"\pi^*"] \arrow[r,"\sim","\pi_{\Gamma'}^*"'] & \mathcal S_2(\Gamma') \arrow[u,"f\mapsto f"']
\end{tikzcd}
\]
\[
\begin{tikzcd}[column sep=large,row sep=large]
\Omega^1(X(\Gamma)) \arrow[r,"\sim","\pi_\Gamma^*"'] \arrow[d,"\pi_*"'] & \mathcal S_2(\Gamma) \arrow[d,"{f\mapsto\sum_{\gamma'\in\Gamma\backslash\Gamma'}f[\gamma']_2}"] \\
\Omega^1(X(\Gamma')) \arrow[r,"\sim","\pi_{\Gamma'}^*"'] & \mathcal S_2(\Gamma')
\end{tikzcd}
\]
\end{lemma*}

\begin{proof}

图中的记号 $\pi_\Gamma^*$ 略有滥用，它是为了提醒我们，同构 $\Omega^1(X(\Gamma))\xrightarrow{\sim}\mathcal S_2(\Gamma)$ 将 $\omega$ 映到满足 $\pi_\Gamma^*(\omega)=f(z)\,dz$ 的 $f$.由此，第一个图的交换性显然成立，因为 $\pi_{\Gamma'}^*=(\pi\circ\pi_\Gamma)^*=\pi_\Gamma^*\circ\pi^*$.



为证明第二个图交换，从 $\omega\in\Omega^1(X(\Gamma))$ 出发，并取 $g\in\mathcal S_2(\Gamma')$ 使 $\pi_{\Gamma'}^*\pi_*\omega=g(z)\,dz$.我们要证明 $g=\sum_{\gamma'\in\Gamma\backslash\Gamma'}f[\gamma']_2$.只须在普通点构成的稠密开集上证明（即避开 $\pi$、$\pi_\Gamma$ 和 $\pi_{\Gamma'}$ 的分歧轨迹）.取普通点 $y\in X(\Gamma')$ 及 $z\in\pi_{\Gamma'}^{-1}(y)$.因为 $\pi_{\Gamma'}$ 在 $y$ 的一个邻域上方是局部同构，可取 $y$ 的邻域 $V_y$ 与 $z$ 的邻域 $O_z$，使 $\pi_{\Gamma'}$ 诱导同构 $O_z\xrightarrow{\sim}V_y$ 且
\begin{align*}
\pi_{\Gamma'}^{-1}(V_y)=\bigsqcup_{\gamma'\in\Gamma'}\gamma'O_z.
\end{align*}
取 $\Gamma\backslash\Gamma'$ 的一个代表元集 $\mathcal L$，便有
\begin{align*}
\pi^{-1}(y)&=\{\pi_\Gamma(\gamma'z),\ \gamma'\in\mathcal L\},\\
\pi^{-1}(V_y)&=\bigsqcup_{\gamma'\in\mathcal L}\pi_\Gamma(\gamma'O_z),
\end{align*}
且对每个 $\gamma'$，$\pi$ 诱导同构 $\pi_\Gamma(\gamma'O_z)\xrightarrow{\sim}V_y$.此外，由 $\pi\circ\pi_\Gamma=\pi_{\Gamma'}$，此同构的逆满足
\begin{align*}
\bigl(\pi|_{\pi_\Gamma(\gamma'O_z)}\bigr)^{-1}\circ\pi_{\Gamma'}=\pi_\Gamma\circ\gamma'.\tag{$*$}\label{eq:modular-2-2-local-inverse}
\end{align*}
现在，由定义有
\begin{align*}
\pi_*(\omega)|_{V_y}=\sum_{\gamma'\in\mathcal L}\Bigl(\bigl(\pi|_{\pi_\Gamma(\gamma'O_z)}\bigr)^{-1}\Bigr)^*\bigl(\omega|_{\pi_\Gamma(\gamma'O_z)}\bigr).
\end{align*}
所以由式~\eqref{eq:modular-2-2-local-inverse}，
\begin{align*}
\bigl(\pi_{\Gamma'}^*(\pi_*\omega)\bigr)|_{O_z}=\sum_{\gamma'\in\mathcal L}\gamma'^*\pi_\Gamma^*\bigl(\omega|_{\pi_\Gamma(\gamma'O_z)}\bigr),
\end{align*}
最终得到
\begin{align*}
(g(z)\,dz)|_{O_z}
&=\sum_{\gamma'\in\mathcal L}\gamma'^*\bigl((f(z)\,dz)|_{\gamma'O_z}\bigr)\\
&=\left(\sum_{\gamma'\in\mathcal L}\gamma'^*(f(z)\,dz)\right)\Big|_{O_z}.
\end{align*}
而 $\gamma'^*(f(z)\,dz)=f(\gamma'z)\,d(\gamma'z)=f[\gamma']_2(z)\,dz$.

\end{proof}



\subsection{对应}\label{sec:modular-2-2-3}

两个紧Riemann曲面 $X$ 与 $X'$ 之间的\newterm{modular-correspondence}{对应}{correspondence} $C$ 是图
\[
\begin{tikzcd}
X & Y \arrow[l,"\pi"'] \arrow[r,"\pi'"] & X'
\end{tikzcd}
\]
其中 $\pi$、$\pi'$ 都是紧Riemann曲面之间的非常值态射.由上一分目，对应 $C$ 诱导同态
\begin{align*}
C_*:=\pi'_*\circ\pi^*:\quad
\left\{\begin{aligned}
\Omega^1(X)&\longrightarrow\Omega^1(X'),\\
\operatorname{Pic}^0(X)&\longrightarrow\operatorname{Pic}^0(X'),\\
H_1(X,\mathbb Z)&\longrightarrow H_1(X',\mathbb Z),\\
J(X)&\longrightarrow J(X'),
\end{aligned}\right.
\end{align*}
它们与 $H_1$ 和 $\Omega^1$ 之间的积分配对及Abel--Jacobi映射相容.



我们关心下述情形.设 $X=X(\Gamma)$，其中 $\Gamma\subset\Gamma(1)$ 是同余子群；设 $\Gamma\alpha\Gamma$ 是 $\Delta$ 中的一个双陪集.与之相伴的是对应 $C_{\Gamma\alpha\Gamma}$：
\[
\begin{tikzcd}[column sep=large]
X(\Gamma) & X(\alpha^{-1}\Gamma\alpha\cap\Gamma) \arrow[l,"\pi^\alpha"'] \arrow[r,"\pi"] & X(\Gamma)
\end{tikzcd}
\]
其中 $\pi^\alpha$ 是复合
\begin{align*}
X(\alpha^{-1}\Gamma\alpha\cap\Gamma)\xrightarrow{\sim}X(\Gamma\cap\alpha\Gamma\alpha^{-1})\longrightarrow X(\Gamma),
\end{align*}
第一个同构由 $\alpha:\mathbb H^*\longrightarrow\mathbb H^*$ 诱导.

\begin{lemma*}

下面的图交换：
\[
\begin{tikzcd}[column sep=large,row sep=large]
\Omega^1(X(\Gamma)) \arrow[r,"\sim","\pi_\Gamma^*"'] \arrow[d,"C_{\Gamma\alpha\Gamma,*}"'] & \mathcal S_2(\Gamma) \arrow[d,"{[\Gamma\alpha\Gamma]_2}"] \\
\Omega^1(X(\Gamma)) \arrow[r,"\sim","\pi_\Gamma^*"'] & \mathcal S_2(\Gamma)
\end{tikzcd}
\]
\end{lemma*}

\begin{proof}

由前一引理，$C_{\Gamma\alpha\Gamma,*}$ 在 $\Omega^1(X(\Gamma))$ 上的作用对应于 $\mathcal S_2(\Gamma)$ 上的下述作用：
\begin{align*}
f\longmapsto\sum_{\gamma\in(\alpha^{-1}\Gamma\alpha\cap\Gamma)\backslash\Gamma}(f[\alpha]_2)[\gamma]_2
=\sum_{\delta\in\Gamma\backslash(\Gamma\alpha\Gamma)}f[\delta]_2
=f[\Gamma\alpha\Gamma]_2,
\end{align*}
即Hecke算子 $[\Gamma\alpha\Gamma]$ 的作用.

\end{proof}

可以直接证明，$C_{\Gamma\alpha\Gamma}$ 的作用在上一分目的各对象 $\operatorname{Pic}^0(X(\Gamma))$、$H_1(X(\Gamma),\mathbb Z)$ 等上定义了Hecke环 $\mathbb Z[\Gamma\backslash\Delta/\Gamma]$ 的作用.不过，现在也可作如下推导：由引理，这在 $\Omega^1$ 上成立；取对偶，便在 $\Omega^{1,*}$ 上成立；由于 $H_1$ 上的作用是 $\Omega^{1,*}$ 上作用的限制，因此在 $H_1$ 上成立；最终便在 $J$ 与 $\operatorname{Pic}^0$ 上成立.

\begin{corollary*}

$\mathbb Z[\Gamma\backslash\Delta/\Gamma]$ 在 $\operatorname{End}_{\mathbb C}(\mathcal S_2(\Gamma))$ 中的像 $\mathbb T_\Gamma$ 是有限生成 $\mathbb Z$ 模.

\end{corollary*}

\begin{proof}

只须证明Hecke环在对偶空间 $\Omega^1(X(\Gamma))^*$ 上的像是有限生成 $\mathbb Z$ 模.



Hecke环保持格 $H_1(X,\mathbb Z)$，所以其像包含于
\begin{align*}
\operatorname{End}_{\mathbb Z}(H_1(X,\mathbb Z)),
\end{align*}
后者是有限生成 $\mathbb Z$ 模.

\end{proof}

现在将讨论限制到 $\Gamma=\Gamma_0(N)$ 及Hecke环 $\mathcal H_0(N)$ 的作用.记其在 $\operatorname{End}_{\mathbb C}(\mathcal S_2(\Gamma_0(N)))$ 中的像为 $\mathbb T_0(N)$；这是一个交换环.



回忆，$T(n)$ 在 $\mathcal S_2(\Gamma_0(N))$ 上的一个特征值系 $(\lambda(n)_{n\in\mathbb N})$ 由素数 $p$ 对应的 $\lambda(p)$ 决定，并确定一个特征（即环同态）$\lambda:\mathcal H_0(N)\longrightarrow\mathbb C$.此特征经 $\mathbb T_0(N)$ 分解；由推论，$\mathbb T_0(N)$ 是有限生成 $\mathbb Z$ 模.因此，$\operatorname{Im}(\lambda)$ 是某个数域 $K_\lambda$ 中的一个整环阶.换言之：

\begin{corollary*}

设 $f\in\mathcal S_2(\Gamma_0(N))$ 对所有 $T(n)$（$n\in\mathbb N$）都是归一化特征形式.则由 $f$ 的 $q$ 展开系数 $a_n(f)$（$n\in\mathbb N$）生成的域 $K_f$ 是数域，且各 $a_n(f)$ 都是其中的代数整数.

\end{corollary*}

为后文之用，需要比较 $\mathcal H_0(N)$ 在 $\mathcal S_2(\Gamma_0(N))$ 中的特征值系与出现在 $H_1(X_0(N),\mathbb C)$ 中的特征值系.用 $V[\lambda]$ 表示 $V$ 中对应于 $\lambda$ 的特征空间，即
\begin{align*}
V[\lambda]=\{v\in V,\ \forall T\in\mathcal H_0(N),\ Tv=\lambda(T)v\}.
\end{align*}
还用 $V_\lambda$ 表示 $V$ 中对应于 $\lambda$ 的广义特征空间，即
\begin{align*}
V_\lambda=\{v\in V,\ \forall T\in\mathcal H_0(N),\ (T-\lambda(T))^nv=0\text{，其中 }n\gg0\}.
\end{align*}
\begin{proposition*}

对任意特征 $\lambda:\mathcal H_0(N)\longrightarrow\mathbb C$，有
\begin{align*}
\dim_{\mathbb C}(H_1(X,\mathbb C)_\lambda)&=2\dim_{\mathbb C}(\mathcal S_2(\Gamma_0(N))_\lambda),\\
\dim_{\mathbb C}(H_1(X,\mathbb C)[\lambda])&>\dim_{\mathbb C}(\mathcal S_2(\Gamma_0(N))[\lambda])\quad\text{（若后者非零）}.
\end{align*}
\end{proposition*}

\begin{proof}

有
\begin{align*}
\dim_{\mathbb C}(\mathcal S_2(\Gamma_0(N))_\lambda)
=\dim_{\mathbb C}(\Omega^1(X_0(N))_\lambda)
=\dim_{\mathbb C}(\Omega^1(X_0(N))_\lambda^*).
\end{align*}
为证明最后一个等式，可取 $\Omega^1(X_0(N))$ 的一组基 $B$，使 $\mathcal H_0(N)$ 在这组基下的作用均为上三角形（对两两交换的自同态族作同时三角化）.于是 $\Omega^1(X_0(N))_\lambda$ 的维数就是 $\lambda$ 在对角线上出现的次数.而在 $\Omega^1(X_0(N))^*$ 的对偶基 $B^*$ 下，$\mathcal H_0(N)$ 由下三角矩阵作用，$\Omega^1(X_0(N))_\lambda^*$ 的维数同样可由对角线读出；这条对角线与基 $B$ 中相同.



另一方面，已知 $H_1(X,\mathbb Z)$ 是 $\Omega^1(X_0(N))^*$ 中的余紧格，所以包含映射
\begin{align*}
H_1(X,\mathbb Z)\hookrightarrow\Omega^1(X_0(N))^*
\end{align*}
诱导 $\mathbb R$ 线性同构
\begin{align*}
H_1(X,\mathbb R)\xrightarrow{\sim}\Omega^1(X_0(N))^*,
\end{align*}
进而诱导 $\mathbb C$ 线性同构
\begin{align*}
H_1(X,\mathbb C)\xrightarrow{\sim}\Omega^1(X_0(N))^*\oplus\overline{\Omega^1(X_0(N))^*},
\end{align*}
其中 $\overline V$ 表示复向量空间的``共轭''，即同一空间，但 $\mathbb C$ 的作用取共轭：对任意 $z\in\mathbb C$ 及 $v\in\overline V$，$z\cdot v:=\overline z v$.当然，$\mathcal H_0(N)$ 仍作用于 $\overline{\Omega^1(X_0(N))^*}$；现在的问题是证明对任意特征 $\lambda$，有
\begin{align*}
\dim_{\mathbb C}(\Omega^1(X_0(N))_\lambda^*)
=\dim_{\mathbb C}\bigl(\overline{\Omega^1(X_0(N))^*}_\lambda\bigr),
\end{align*}
也就是
\begin{align*}
\dim_{\mathbb C}(\mathcal S_2(\Gamma_0(N))_\lambda)
=\dim_{\mathbb C}\bigl(\overline{\mathcal S_2(\Gamma_0(N))^*}_\lambda\bigr).
\end{align*}
注意到Petersson内积通过 $g\mapsto(f\mapsto\langle f,g\rangle)$ 给出同构 $\mathcal S_2(\Gamma_0(N))\xrightarrow{\sim}\overline{\mathcal S_2(\Gamma_0(N))^*}$.但各 $T(n)$ 对此内积并不全都自伴（只有 $(n,N)=1$ 时才自伴），所以此同构与 $\mathcal H_0(N)$ 的作用并非``完全''相容.不过，我们也已看到，当 $p\mid N$ 时，$[T(p)]_2$ 的伴随为
\begin{align*}
[T(p)]_2^*=[w_N]_2\circ[T(p)]_2\circ[w_N^{-1}]_2.
\end{align*}
此式事实上对所有素数 $p$ 都成立，因为当 $p\nmid N$ 时，$[T(p)]_2$ 与 $[w_N]_2$ 交换.由此，同构
\begin{align*}
\mathcal S_2(\Gamma_0(N))&\xrightarrow{\sim}\overline{\mathcal S_2(\Gamma_0(N))^*},\\
g&\longmapsto\varphi_g:\bigl(f\longmapsto\langle f[w_N^{-1}]_2,g\rangle\bigr)
\end{align*}
与 $\mathcal H_0(N)$ 的作用相容（即 $\varphi_{g[T]_2}(f)=\varphi_g(f[T]_2)$），因此对任意特征 $\lambda:\mathcal H_0(N)\longrightarrow\mathbb C$，有
\begin{align*}
\dim_{\mathbb C}(\mathcal S_2(\Gamma_0(N))_\lambda)
&=\dim_{\mathbb C}\bigl(\overline{\mathcal S_2(\Gamma_0(N))^*}_\lambda\bigr),\\
\dim_{\mathbb C}(\mathcal S_2(\Gamma_0(N))[\lambda])
&=\dim_{\mathbb C}\bigl(\overline{\mathcal S_2(\Gamma_0(N))^*}[\lambda]\bigr).
\end{align*}
这就证明了所述等式.不等式也随之成立，因为若
\begin{align*}
\dim(\Omega^1(X_0(N))[\lambda])\ne0,
\end{align*}
则 $\dim(\Omega^1(X_0(N))^*_{\lambda})\ne0$，因而 $\dim(\Omega^1(X_0(N))^*[\lambda])\ne0$.

\end{proof}

\begin{corollary*}

设 $f\in\mathcal S_2(\Gamma_0(N))$ 是本原形式，$\lambda$ 为相应的特征值系.则
\begin{align*}
\dim_{\mathbb C}\bigl(H_1(X_0(N),\mathbb C)[\lambda]\bigr)
=\dim_{K_f}\bigl(H_1(X_0(N),K_f)[\lambda]\bigr)=2.
\end{align*}
\end{corollary*}

\begin{proof}

已知此时 $\mathcal S_2(\Gamma_0(N))[\lambda]=\mathcal S_2(\Gamma_0(N))_\lambda$ 是一维的.由命题，
\begin{align*}
\dim_{\mathbb C}\bigl(H_1(X_0(N),\mathbb C)[\lambda]\bigr)
&\leqslant\dim_{\mathbb C}\bigl(H_1(X_0(N),\mathbb C)_\lambda\bigr)=2,
\end{align*}
而且 $\dim_{\mathbb C}(H_1(X_0(N),\mathbb C)[\lambda])>1$.所以 $\dim_{\mathbb C}(H_1(X_0(N),\mathbb C)[\lambda])=2$.由于 $\lambda$ 的值属于子域 $K_f\subset\mathbb C$，也有 $\dim_{K_f}(H_1(X_0(N),K_f)[\lambda])=2$.

\end{proof}

继续取推论中的 $f$.现在设 $\ell$ 是素数，$\mathfrak l\in\operatorname{Max}(\mathcal O_{K_f})$ 是位于 $\ell$ 上方的极大理想.置
\begin{align*}
V_f:=H_1(X,K_{f,\mathfrak l})[\lambda].
\end{align*}
这个 $K_{f,\mathfrak l}$ 上的二维空间就是我们所寻求的Galois表示的空间.为定义此表示，将使用它与 $\operatorname{Pic}^0$ 的下述联系：
\begin{align*}
H_1(X,K_{f,\mathfrak l})
=\left(\varprojlim_m\operatorname{Pic}^0(X_0(N))[\ell^m]\right)\otimes_{\mathbb Z_\ell}K_{f,\mathfrak l}.
\end{align*}
\begingroup\emergencystretch=3em
事实上，若以 $\operatorname{Pic}^0(X)[N]$ 表示 $\operatorname{Pic}^0$ 上乘以 $N$ 的映射之核，则Abel--Jacobi定理将 $\operatorname{Pic}^0(X)[N]$ 与
\par\endgroup
\begin{align*}
N^{-1}H_1(X,\mathbb Z)/H_1(X,\mathbb Z)\simeq H_1(X,\mathbb Z)\otimes(\mathbb Z/N\mathbb Z)
\end{align*}
等同；取投射极限便得到同构
\begin{align*}
\varprojlim_m\operatorname{Pic}^0(X)[\ell^m]
\simeq H_1(X,\mathbb Z)\otimes\mathbb Z_\ell=H_1(X,\mathbb Z_\ell).
\end{align*}
由此可见，现在需要在 $\operatorname{Pic}^0(X_0(N))$ 上定义Galois作用.为此，必须先为 $X_0(N)$ 寻找一个作为 $\mathbb Q$ 上代数曲线的``典范''模型.

